\textbf{Operaciones para comparar ambos casos.}
\begin{enumerate}

\item Despejar empuje de potencia cinética: F=2P/v.

\[\fitmath{\mathrm{FA}=\frac{2\cdot 600}{3000}\approx 0.4\,\mathrm{N}}\]

\item Del balance de energía, caudal es 2P/v al cuadrado.

\[\fitmath{\mathrm{mA}=\frac{2\cdot 600}{{3000}^{2}}\approx 0.0001333333\,\mathrm{kg/s}}\]

\item Integrar empuje constante durante el intervalo.

\[\fitmath{\mathrm{IA}=\mathrm{FA}\cdot 120}\]

\[\fitmath{\mathrm{IA}=(0.4)\cdot 120\approx 48\,\mathrm{N\,s}}\]

\item Repetir el balance de momento y energía con la segunda velocidad.

\[\fitmath{\mathrm{FB}=\frac{2\cdot 600}{18000}\approx 0.06666667\,\mathrm{N}}\]

\item Calcular el segundo caudal.

\[\fitmath{\mathrm{mB}=\frac{2\cdot 600}{{18000}^{2}}\approx 3.703704\times10^{-6}\,\mathrm{kg/s}}\]

\item Impulso del segundo caso.

\[\fitmath{\mathrm{IB}=\mathrm{FB}\cdot 120}\]

\[\fitmath{\mathrm{IB}=(0.06666666667)\cdot 120\approx 8\,\mathrm{N\,s}}\]

\end{enumerate}
